\section{Methodology}
\label{sec:methodology}

This section presents GMM-Router, a framework for parameter-efficient continual code generation that routes each query to task-specific experts without retraining earlier adapters or retaining raw training data. GMM-Router has four components: task-specific LoRA expert training (Section~\ref{sec:experts}); routing feature extraction through a frozen backbone and a fixed orthonormal projection (Section~\ref{sec:features}); a continual GMM instruction router that supports hard and soft routing (Sections~\ref{sec:router} and~\ref{sec:soft_routing}); and, for the executable benchmark, an iterative refinement pipeline built around a second, traceback-based router (Section~\ref{sec:traceback}).

At a high level, GMM-Router processes the task sequence as follows. For each new task, a dedicated LoRA adapter and a GMM instruction router are trained in parallel (Sections~\ref{sec:experts}--\ref{sec:router}). At inference, the instruction router scores the query against every trained task and produces a routing distribution over adapters, which is used to select the highest-weighted adapter for hard routing or to merge all adapters for soft routing (Section~\ref{sec:soft_routing}); this constitutes Round-1 inference and is the only inference mechanism used on the CodeTask benchmark. On the executable benchmark, adapters are additionally evaluated on a held-out calibration split (Section~\ref{sec:setup}) whose failed predictions are used to refine each adapter on repair pairs and to fit a traceback GMM on error messages (Section~\ref{sec:traceback}); for any test sample whose five Round-1 candidates all fail execution, a second inference round combines the instruction and traceback routing signals through a Jensen--Shannon-divergence (JSD) gated rule to produce one additional prediction (Section~\ref{sec:round2}).

% ---------------------------------------------------------------------------
\subsection{Training task-specific experts}
\label{sec:experts}
% ---------------------------------------------------------------------------

For each task $\mathcal{T}_k$ with training dataset $\mathcal{D}_k$ (Eq.~\eqref{eq:dataset}), GMM-Router trains a task-specific LoRA adapter \citep{hu2022lora} on top of a shared frozen backbone, following the LoRA formulation in Section~\ref{sec:peft_background}. Each adapter is a low-rank update applied to the query and value projections in every transformer layer. Only the adapter parameters $\Theta_k=\{B_k^{(\ell)},A_k^{(\ell)}\}_\ell$ are updated; the backbone and all previously trained adapters remain frozen throughout training, so a new task cannot interfere with previously learned ones. The training objective is the negative log-likelihood of the target sequence,
\begin{equation}
    \mathcal{L}_k(\Theta_k) = -\sum_{(x,y)\in\mathcal{D}_k} \sum_{t=1}^{|y|}
    \log p_k(y_t \mid x, y_{<t}),
    \label{eq:task_loss}
\end{equation}
where $p_k(y_t\mid x,y_{<t})$ is the output distribution of expert $k$ at decoding step $t$. The backbone is split at layer $\ell^*$: layers $0$ through $\ell^*{-}1$ produce the frozen features used for routing (Section~\ref{sec:features}), while LoRA adapters modify layers $\ell^*$ through $L$. A larger $\ell^*$ dedicates more layers to routing features but leaves fewer layers for task-specific adaptation.

% ---------------------------------------------------------------------------
\subsection{Routing feature extraction}
\label{sec:features}
% ---------------------------------------------------------------------------

Stable routing requires a fixed representation space: if the feature map changed as new tasks arrived, GMMs fitted on earlier tasks would become misaligned with those fitted later. GMM-Router achieves this by extracting features from the frozen part of the backbone --- with all LoRA adapters disabled, so the routing space is independent of any per-task adapter weights --- and projecting them with a random matrix drawn once and fixed for the entire continual training run.

For a tokenized input of length $T$, the backbone produces hidden states $\mathbf{H}^{(\ell)}\in\mathbb{R}^{T\times d}$ at each layer $\ell$; the routing representation is taken at layer $\ell^*$, $\mathbf{H}=\mathbf{H}^{(\ell^*)}$. Using an intermediate layer yields representations that are less task-specific and more stable as tasks accumulate. The sequence-level representation is the mean of the token hidden states,
\begin{equation}
    \tilde{\mathbf{h}} = \frac{1}{T}\sum_{t=1}^{T} \mathbf{h}_t \in \mathbb{R}^d,
    \label{eq:meanpool}
\end{equation}
and parameter-free layer normalization removes scale differences,
\begin{equation}
    \mathbf{h} = \operatorname{LN}(\tilde{\mathbf{h}}), \qquad
    \operatorname{LN}(\mathbf{v}) = \frac{\mathbf{v} - \mu_{\mathbf{v}}}{\sigma_{\mathbf{v}} + \varepsilon},
    \label{eq:ln}
\end{equation}
where $\mu_\mathbf{v}$ and $\sigma_\mathbf{v}$ are the mean and standard deviation of the elements of $\mathbf{v}$. A fixed projection matrix $\mathbf{P}\in\mathbb{R}^{p\times d}$ then reduces the dimension to $p\ll d$: a random Gaussian matrix is drawn with a fixed seed, its thin QR factorization gives $\mathbf{P}$, so that $\mathbf{P}\mathbf{P}^\top=\mathbf{I}_p$. The routing embedding for input $x$ is
\begin{equation}
    \mathbf{z} = \mathbf{P}\,\mathbf{h} \in \mathbb{R}^p.
    \label{eq:z}
\end{equation}
Because $\mathbf{P}$ never changes, the routing embedding (Eq.~\eqref{eq:z}) lies in a fixed $p$-dimensional space for every task, so GMM densities fitted at any point in the sequence remain directly comparable to those fitted later.

% ---------------------------------------------------------------------------
\subsection{Continual GMM instruction router}
\label{sec:router}
% ---------------------------------------------------------------------------

Each task's instruction distribution is modeled as a diagonal GMM with $M$ components (Eq.~\eqref{eq:gmm}), fitted by EM (Algorithm~\ref{alg:em_gmm}) on the routing embeddings $\{\mathbf{z}_i^{(k)}\}_{i=1}^N$ of task $k$. When task $k$ arrives, its GMM is fitted and appended to the router registry; no parameters belonging to earlier tasks are modified.

\textbf{MAP scoring with a uniform class prior.} At inference, the router scores each registered task $k$ using the log-posterior under a uniform prior over the $K$ trained tasks,
\begin{equation}
    s_k(\mathbf{z}) = \log p_k(\mathbf{z}) + \log \frac{1}{K}.
    \label{eq:snorm}
\end{equation}
Since $\log(1/K)$ is constant across tasks, the argmax decision is equivalent to selecting the highest log-likelihood; the prior term keeps scores on a consistent scale as $K$ grows.

\textbf{Hard routing.} Given a query embedding $\mathbf{z}$, the hard router selects the task with the highest score,
\begin{equation}
    \hat{k} = \operatorname*{argmax}_{k \in \{1,\ldots,K\}} \; s_k(\mathbf{z}),
    \label{eq:routing}
\end{equation}
and generation proceeds using only expert $\hat{k}$'s adapter. This is the default inference strategy for the CodeTask benchmark.

% ---------------------------------------------------------------------------
\subsection{Soft routing and adapter merging}
\label{sec:soft_routing}
% ---------------------------------------------------------------------------

Rather than committing to a single expert, soft routing converts the log-posterior scores into a probability distribution over all $K$ tasks via a temperature-scaled softmax,
\begin{equation}
    \alpha_k(x) = \frac{\exp\!\left(s_k(\mathbf{z}) / \tau\right)}{\sum_{j=1}^K \exp\!\left(s_j(\mathbf{z}) / \tau\right)}, \qquad \bm{\alpha}(x) \in \Delta^{K-1},
    \label{eq:alpha}
\end{equation}
where $\tau>0$ is the routing temperature; as $\tau\to0^+$, $\bm{\alpha}(x)\to\mathbf{e}_{\hat{k}}$ recovers hard routing, and as $\tau\to\infty$, $\bm{\alpha}(x)\to\mathbf{1}/K$ becomes uniform. The $K$ LoRA adapters are then merged at the weight level,
\begin{equation}
    \Delta W_{\mathrm{merged}} = \sum_{k=1}^K \alpha_k(x)\, B_k A_k,
    \label{eq:weight_merge}
\end{equation}
yielding a single merged adapter applied to the frozen backbone for generation, at the cost of only one forward pass per decoding step regardless of $K$. Weight-level merging (Eq.~\eqref{eq:weight_merge}) approximates the mixture-of-experts output distribution $p_{\mathrm{mix}}(y_t\mid x,y_{<t})=\sum_k \alpha_k(x)\,p_k(y_t\mid x,y_{<t})$, which would otherwise require $K$ independent forward passes; because $\alpha_k(x)$ depends only on $x$, the routing weights are computed once per query and held fixed throughout autoregressive decoding. Algorithm~\ref{alg:generation} summarizes soft-routing generation, and Figure~\ref{fig:arch} illustrates the resulting Round-1 inference pipeline.

\begin{algorithm}[t]
\caption{Soft-Routing Expert-Mixture Generation}
\label{alg:generation}
\KwIn{Query $x$; log-posterior scores $\{s_k(\mathbf{z})\}$; adapters $\{(B_k, A_k)\}_{k=1}^K$; temperature $\tau$}
\KwOut{Generated output $\hat{y}$}
Extract routing embedding $\mathbf{z}$ with LoRA adapters disabled\;
$\alpha_k \leftarrow \exp(s_k(\mathbf{z})/\tau)/\sum_j\exp(s_j(\mathbf{z})/\tau)$, $\forall k$\;
Merge adapters: $\Delta W \leftarrow \sum_k \alpha_k B_k A_k$\;
Generate $\hat{y}$ with the merged adapter\;
\Return $\hat{y}$\;
\end{algorithm}

\begin{figure}[t]
\centering
\resizebox{\linewidth}{!}{%
\begin{tikzpicture}[
    font=\small,
    shared/.style={draw=blue!55!black, line width=0.8pt, rounded corners=5pt,
        minimum width=2.2cm, minimum height=0.90cm, align=center, fill=blue!10},
    znode/.style={draw=blue!55!black, line width=0.8pt, rounded corners=5pt,
        minimum width=1.0cm, minimum height=0.90cm, align=center, fill=blue!20},
    rnode/.style={draw=violet!65!black, line width=0.8pt, rounded corners=5pt,
        minimum width=2.6cm, minimum height=0.90cm, align=center, fill=violet!10},
    enode/.style={draw=teal!65!black, line width=0.8pt, rounded corners=5pt,
        minimum width=1.9cm, minimum height=0.90cm, align=center, fill=teal!10},
    outnode/.style={draw=orange!65!black, line width=0.8pt, rounded corners=5pt,
        minimum width=1.9cm, minimum height=0.90cm, align=center, fill=orange!12},
    arr/.style= {-{Latex[length=5pt,width=4pt]}, thick, draw=gray!55},
    rarr/.style={-{Latex[length=5pt,width=4pt]}, thick, draw=violet!60},
    earr/.style={-{Latex[length=5pt,width=4pt]}, thick, draw=teal!60},
    oarr/.style={-{Latex[length=5pt,width=4pt]}, thick, draw=orange!70!black},
]
\node[shared]                          at ( 0.0, 0) (x)      {Input $x$};
\node[shared]                          at ( 2.7, 0) (enc)    {Backbone\\layers $0$--$\ell^*$};
\node[znode]                           at ( 5.0, 0) (z)      {$\mathbf{z}$};
\node[rnode]                           at ( 7.4, 0) (gmm)    {GMM Router\\$\log p_k(\mathbf{z})$};
\node[rnode, minimum width=2.4cm]      at (10.5, 0) (scores) {Scores $s_k(\mathbf{z})$};
\node[rnode, minimum width=2.0cm]      at (13.4, 0) (alpha)  {$\bm{\alpha}(x)$};
\draw[arr]  (x)      -- (enc);
\draw[arr]  (enc)    -- (z);
\draw[rarr] (z)      -- (gmm);
\draw[rarr] (gmm)    -- (scores);
\draw[rarr] (scores) -- (alpha);
\node[enode, minimum width=4.0cm] at (7.8, -2.8) (merge) {Merge LoRA:\\$\Delta W = \sum_k \alpha_k B_k A_k$};
\node[outnode] at (12.5, -2.8) (out) {Generate\\$\hat{y}$};
\draw[earr] (z.south) to[out=-90, in=180] (merge.west);
\draw[rarr] (alpha.south) to[out=-145, in=90] (merge.north);
\draw[oarr] (merge) -- (out);
\begin{scope}[on background layer]
    \fill[gray!6, rounded corners=10pt]
        ($(x.north west)+(-0.18, 0.47)$) rectangle ($(out.south east)+(1.36,-0.42)$);
    \node[draw=violet!28, fill=violet!5, rounded corners=9pt, inner sep=6pt,
          fit=(gmm)(scores)(alpha)] (rpanel) {};
    \node[draw=teal!28, fill=teal!5, rounded corners=9pt, inner sep=6pt,
          fit=(merge)] (epanel) {};
\end{scope}
\node[font=\scriptsize\itshape, violet!70] at (rpanel.north) [above=2pt] {Router branch};
\node[font=\scriptsize\itshape, teal!70]   at (epanel.south) [below=2pt] {Expert branch};
\end{tikzpicture}}
\caption{Round-1 inference pipeline of GMM-Router. The frozen backbone maps input $x$ to routing embedding $\mathbf{z}$ using layers $0$--$\ell^*$. The router branch (purple) evaluates the $K$ GMM log-densities, adds the uniform log-prior, and converts scores to routing weights $\bm{\alpha}(x)$. The expert branch (teal) merges the $K$ LoRA adapters as $\Delta W=\sum_k\alpha_k B_k A_k$ and generates output $\hat{y}$.}
\label{fig:arch}
\end{figure}

% ---------------------------------------------------------------------------
\subsection{Iterative refinement for executable code generation}
\label{sec:traceback}
% ---------------------------------------------------------------------------

For executable code generation, GMM-Router supports a second refinement round for samples where all Round-1 predictions fail execution. A separate GMM traceback router is trained on runtime error messages rather than instruction prompts, providing a complementary routing signal, and each per-task adapter is optionally refined on repair pairs; both procedures use the calibration predictions described in Section~\ref{sec:setup}.

\textbf{Traceback router.} The traceback router is a diagonal GMM (Section~\ref{sec:gmm_background}) fitted on error-message embeddings extracted from failed calibration predictions. For each failed candidate with a non-empty error string, the error text is collected and deduplicated; the resulting strings are encoded with the same frozen backbone used for instruction routing and projected into a separate routing space with its own fixed projection matrix, and a GMM is fitted per task using the EM procedure of Section~\ref{sec:gmm_background}.

\textbf{Adapter refinement.} In parallel with traceback router training, each per-task LoRA adapter can be further fine-tuned on repair pairs derived from the evaluated calibration data. Each pair consists of the original instruction, buggy code, and runtime traceback as input, paired with a ground-truth or passing candidate as target: for each calibration sample, the first failed candidate with a non-empty error string supplies the (buggy code, traceback) pair, and the repair target is the first passing candidate, or the ground-truth solution if none of the five candidates passed; samples with no failed candidate or no repair target are skipped.

\textbf{Round-2 inference and JSD-gated score combination.}
\label{sec:round2}
Round-2 inference applies only to samples where every Round-1 candidate failed. Let $\mathbf{s}_{\mathrm{in}}\in\mathbb{R}^K$ be the instruction router's Round-1 log-posterior scores and $\mathbf{s}_{\mathrm{tb}}\in\mathbb{R}^K$ the traceback router's scores on the corresponding error message; each is converted to a distribution, $\bm{\alpha}_{\mathrm{in}}=\operatorname{softmax}(\mathbf{s}_{\mathrm{in}}/\tau)$ and $\bm{\alpha}_{\mathrm{tb}}=\operatorname{softmax}(\mathbf{s}_{\mathrm{tb}})$. The two routers use fundamentally different signals --- natural-language problem descriptions versus runtime error messages --- so their agreement is informative about routing confidence: when both assign high probability to the same expert, that consensus is strong evidence the expert is correct; when they disagree, neither signal is fully trustworthy in isolation, and committing to either alone risks compounding the original routing error.

This motivates gating the combination by the degree of conflict between the two distributions. Let $\mathbf{m}=\tfrac{1}{2}(\bm{\alpha}_{\mathrm{in}}+\bm{\alpha}_{\mathrm{tb}})$. The Kullback--Leibler divergence from $p$ to $q$ is $\mathrm{KL}(p\|q)=\sum_k p_k\log(p_k/q_k)$ \citep{kullback1951information}, and the Jensen--Shannon divergence \citep{lin1991divergence} is its symmetrized, bounded variant,
\begin{equation}
    \mathrm{JSD}(\bm{\alpha}_{\mathrm{in}} \| \bm{\alpha}_{\mathrm{tb}})
    = \tfrac{1}{2}\mathrm{KL}(\bm{\alpha}_{\mathrm{in}}\|\mathbf{m})
    + \tfrac{1}{2}\mathrm{KL}(\bm{\alpha}_{\mathrm{tb}}\|\mathbf{m}).
    \label{eq:jsd}
\end{equation}
The combined routing weight is
\begin{equation}
    \bm{\alpha} = (1 - \mathrm{JSD})\operatorname{softmax}(\mathbf{s}_{\mathrm{in}} + \mathbf{s}_{\mathrm{tb}}) + \mathrm{JSD} \cdot \frac{\mathbf{1}}{K}.
    \label{eq:jsd_gated}
\end{equation}
When the two routers agree (low JSD), the combined log-posterior $\mathbf{s}_{\mathrm{in}}+\mathbf{s}_{\mathrm{tb}}$ is trusted and converted into a peaked distribution; when they disagree strongly (high JSD), the distribution is blended toward uniform, spreading probability mass across all experts rather than committing to a potentially wrong one. The rule in Eq.~\eqref{eq:jsd_gated} requires no additional learned parameters or tuned thresholds. The final Round-2 prediction is generated using $\bm{\alpha}$ and the same weight-level adapter merging as in Round-1 (Eq.~\eqref{eq:weight_merge}).

In summary, task-specific LoRA adapters prevent catastrophic forgetting through architectural isolation, since training on a new task updates only its dedicated adapter while all earlier adapters and the backbone remain frozen. Routing feature extraction projects frozen backbone representations through a fixed orthonormal matrix, keeping the routing space stable throughout the task sequence so GMM densities fitted at any point remain comparable as more tasks arrive. A per-task diagonal GMM fitted by EM models each task's instruction distribution and provides log-posterior scores for task-agnostic routing, realized as hard argmax selection for the CodeTask benchmark or temperature-scaled soft adapter merging for the executable benchmark; for the latter, the iterative refinement pipeline recovers samples that exhaust all Round-1 candidates by combining instruction and traceback routing signals through the JSD-gated rule.
